\section{Анализ производительности верификационного стенда}

Оценка вычислительной эффективности предложенного подхода проводилась при моделировании полного цикла работы навигационного приемника. Исходный массив оцифрованного сигнала имеет длительность 8 секунд, что необходимо для последующей работы канала слежения. Однако блок первичного поиска анализирует только фрагмент длительностью 1 мс, накапливая 4096 отсчетов после децимации. Поскольку алгоритм PCPS проверяет каждую из 201 доплеровской гипотезы независимо, симулятору требуется выполнить полный цикл конвейерной обработки (БПФ, перемножение и ОБПФ) для каждого среза. В результате полный цикл симуляции блока поиска потребовал непрерывного расчета \(11{,}87 \cdot 10^6\) тактов тактовой частоты схемы. После успешного обнаружения сигнала блок поиска отключается, а управление передается каналу слежения.

Ключевые параметры производительности среды на базе Verilator для этапа обнаружения приведены в табл.~\ref{tab:perf}. Пропускная способность симуляции достигла \(0{,}150\)~млн тактов в секунду, а общее время выполнения процесса составило около 80 секунд. 

Столь значительное время выполнения обусловлено высокой вычислительной сложностью алгоритма PCPS: для обработки всего одной миллисекунды сигнала требуется выполнить 201 независимый цикл согласованной фильтрации (каждый включает прямое БПФ, комплексное перемножение, обратное БПФ и квадратичное детектирование). Суммарно симулятору необходимо рассчитать более 400 преобразований Фурье размерностью 4096 точек, что эквивалентно \(11{,}87 \cdot 10^6\) тактов работы синхронного конвейера. При этом в отличие от классических событийно-управляемых симуляторов (Vivado XSIM, ModelSim), где на каждом такте формируются миллионы событий переключения триггеров и требуется их последовательная обработка через очередь планировщика, архитектура Verilator выполняет модель как единый нативный C++ процесс. 

Полученный показатель времени симуляции подтверждает пригодность данного подхода для практической верификации. Полный перебор 201 частотного среза для 1 мс выборки занимает менее полутора минут машинного времени. Это делает возможной итеративную отладку алгоритмов ЦОС, тогда как аналогичный тест в классических САПР требует значительно большего времени и вычислительных ресурсов.


\begin{table}[htbp]
\caption{Coarse Acquisition Performance}
\begin{center}
\begin{tabular}{|l|c|}
\hline
\textbf{Metric} & \textbf{Value} \\
\hline
Doppler bins & 201 \\
\hline
Simulated input-clock cycles & 11867040 \\
\hline
Acquisition time & 79.227 s \\
\hline
Total process time & 80.637 s \\
\hline
Simulation throughput & 0.150 Mcycles/s \\
\hline
Equivalent input time at 3.819 Msamples/s & 3.107 s \\
\hline
Speedup vs. input real time & 0.039 \\
\hline
\end{tabular}
\label{tab:perf}
\end{center}
\end{table}
